\section*{Заключение}
\addcontentsline{toc}{section}{Заключение}
\thispagestyle{plain}

% TODO: ~3 страницы.

В рамках настоящей магистерской работы решена задача разработки
моделей и сервиса для интерактивного взаимодействия с обучающимися
в условиях современного высшего учебного заведения.

\textbf{Основные результаты работы:}

\begin{enumerate}[label=\arabic*., noitemsep]
    \item Проведён анализ современных систем интерактивного обучения
          и подходов к геймификации; выявлены ограничения предшествующей
          научно-исследовательской работы и обоснован переход
          к собственному backend-решению, развёрнутому в инфраструктуре
          Российской Федерации.
    \item Разработаны модели взаимодействия с обучающимися:
          ролевая модель, модель курса и банка вопросов, модель
          интерактивной сессии, модель игровой прогрессии, модель
          интервального повторения SM\textendash 2 и модель оценки
          надёжности усвоения знаний.
    \item Спроектирована и реализована клиент-серверная архитектура
          сервиса на стеке Flutter + Go + PostgreSQL + Redis + WebSocket;
          реализована интеграция с университетской LMS edu.gubkin.
          Клиент собирается под платформы Web, iOS, Android и macOS.
    \item Выполнен расчёт показателей надёжности архитектуры
          (вероятности безотказной работы, средней наработки на отказ,
          коэффициента готовности); реализованы резервирование PostgreSQL,
          мониторинг и нагрузочное тестирование. Тем самым теория
          надёжности, преподаваемая на кафедре АСУ, применена
          к собственной архитектуре сервиса.
    \item Проведён педагогический эксперимент со студентами курса
          <<Надёжность>>; получены данные о влиянии разработанного
          сервиса и геймификации на успеваемость и удержание знаний.
\end{enumerate}

\textbf{Научная новизна} результатов состоит в построении универсального
сервиса интерактивного тестирования с двусторонней интеграцией с LMS
edu.gubkin, в применении расширенной геймификации и интервального
повторения поверх стандартных тестовых форматов, и в формировании
замкнутой инженерно-педагогической петли <<обучение надёжности
на сервисе с расчётно обоснованной надёжностью>>.

\textbf{Практическая значимость.}
Сервис развёрнут в инфраструктуре Российской Федерации и готов
к использованию в образовательном процессе РГУ Нефти и Газа
им. И.М. Губкина для дисциплин любых кафедр.

\textbf{Направления дальнейшего развития:}
расширение библиотеки интерактивных инструментов; адаптивная
геймификация на основе таксономии Бартла; внедрение Cloud Functions
для критичной бизнес-логики; интеграция с другими университетскими
системами; распространение на иные образовательные учреждения.

\newpage
